\documentclass[aspectratio=169]{beamer}
\usetheme{Madrid}
\usecolortheme{seahorse}
\usepackage{hyperref}
\usepackage{amsmath,amssymb}

\title{A3 + The Slop Feedback Loop}
\subtitle{Formal Methods Deep Dive (Private Notes)}
\author{Halley Young}
\date{March 2026}

\setbeamertemplate{navigation symbols}{}
\setbeamertemplate{footline}[frame number]

\newcommand{\syncline}[2]{\textbf{#1} \hfill \textcolor{gray}{#2}}

\begin{document}

\begin{frame}
\titlepage
\end{frame}

\begin{frame}{1. Deliverable first}
\syncline{Slide 1}{~2 min}
\textbf{Script:}
\begin{itemize}
  \item Open with what ships, not with history: \texttt{pip install a3-python} gives scan + static elimination + optional triage.
  \item The architecture is explicitly static-first, agentic-second.
  \item Stress asymmetry: formal methods take the bulk volume; LLM handles uncertainty residue only.
  \item Set expectation for RISE audience: this talk is about what is proved, what is approximated, and where assumptions live.
  \item Keep phrase: ``proof/disproof artifacts are upstream; triage is downstream policy support.''
\end{itemize}
\end{frame}

\begin{frame}{2. Blog arc we follow}
\syncline{Slide 2}{~2 min}
\textbf{Script:}
\begin{itemize}
  \item Follow the blog storyline directly: origin challenge \(\rightarrow\) Episode I/II/III \(\rightarrow\) shipping architecture.
  \item Then expand on RISE-sensitive details: transition semantics, obligations, solver outcomes, refinement policy.
  \item End on operational policy: kitchensink orchestration, residual triage, CI ratchet.
  \item Frame from the start: this is verification-oriented systems engineering, not theorem theater.
\end{itemize}
\end{frame}

\begin{frame}{3. One ridiculous origin story}
\syncline{Slide 3}{~2 min}
\textbf{Script:}
\begin{itemize}
  \item Blog claim: this started as ``can AI generate paper-scale verification ideas that survive code?''
  \item Early theory looked elegant but collapsed on implementation details.
  \item Working rule became adversarial: every idea must survive tests and regression runs.
  \item Useful phrase from blog: ``fight AI slop with AI slop'' \(\rightarrow\) generate aggressively, filter brutally.
  \item This is why the stack values executable semantics above conceptual novelty.
\end{itemize}
\end{frame}

\begin{frame}{4. Episode I: quantitative safety intuition}
\syncline{Slide 4}{~2 min}
\textbf{Script:}
\begin{itemize}
  \item Distance semantics survived as prioritization intuition, not final correctness engine.
  \item Keep distinction explicit: ranking answers ``what first,'' verification answers ``what never.''
  \item This sets up Episode II naturally: separation proofs for reachability exclusion.
  \item If challenged, concede limits of pure metrics for industrial false-positive suppression.
\end{itemize}
\end{frame}

\begin{frame}{5. Episode II: prove unreachable}
\syncline{Slide 5}{~2 min}
\textbf{Script:}
\begin{itemize}
  \item Reachability framing: show \(R\cap U=\varnothing\) under encoded semantics and assumptions.
  \item Barrier-style certificates matter because they summarize many executions with inductive constraints.
  \item They are only meaningful when tied to concrete state projections from actual Python execution points.
  \item Failure is informative: obligation counterexamples feed refinement and portfolio rerouting.
  \item Good caveat: certificates are powerful but not universally synthesizeable under budget.
\end{itemize}
\end{frame}

\begin{frame}{6. Barrier obligations (informal form)}
\syncline{Slide 6}{~2 min}
\textbf{Script:}
\begin{itemize}
  \item Initialization: \(\forall s\in S_0,\;B(s)\ge\epsilon\).
  \item Consecution: \(\forall s\to s',\;B(s)\ge0\Rightarrow B(s')\ge0\).
  \item Unsafe separation: \(\forall s\in U,\;B(s)\le-\epsilon\).
  \item Emphasize conjunctive acceptance: one failed obligation means no certificate.
  \item Solver policy is conservative: UNKNOWN does not justify pruning.
  \item RISE point: this is a soundness-preserving operational policy, not convenience behavior.
\end{itemize}
\end{frame}

\begin{frame}[t]{6b. Exact workflow: one polynomial certificate, one Python bug}
\syncline{Slide 6b}{~3 min}
\textbf{Script:}
\scriptsize
\begin{itemize}
  \item \textbf{Bug anchor:} pick one reported site, e.g., DIV\_ZERO at \texttt{f} offset \(o\) for \texttt{n / d}.
  \item \textbf{State basis:} choose symbolic variables \(x\) from locals/stack/heap observers at \(o\); compute path precondition \(\Phi(x)\).
  \item \textbf{Unsafe set:} \(U(x)\equiv(d(x)=0)\land\Phi(x)\).
  \item \textbf{Dynamics:} derive entry set \(S_0\) and transition relation \(T(x,x')\) from bytecode transfer semantics.
  \item \textbf{Template:} choose degree-\(k\) polynomial \(B(x)=\sum_{\alpha\in\mathcal{M}_k}c_\alpha x^\alpha\) with unknown coefficients \(c_\alpha\).
  \item \textbf{Obligations:} init \(\forall x\in S_0: B(x)\ge\epsilon\); consecution \(\forall x,x': T\land B(x)\ge0\Rightarrow B(x')\ge0\); unsafe \(\forall x: U(x)\Rightarrow B(x)\le-\epsilon\).
  \item \textbf{Solve:} synthesize \(c_\alpha\) via staged backends (fast checks first, stronger relaxations next).
  \item \textbf{Validate:} replay obligations on current model; any counterexample rejects candidate certificate.
  \item \textbf{Mini-example:} if Python guard is \texttt{if den > 0: return num/den}, choose \(x=den\), \(U(x)\equiv(x=0)\), and template \(B(x)=x-\tfrac{1}{2}\). Then init/consecution hold under guard-preserving transitions and unsafe separation holds since \(B(0)=-\tfrac{1}{2}<0\).
  \item \textbf{Map back:} valid certificate marks this exact site unreachable (FP) with proof artifact: template, coefficients, obligations, site metadata.
  \item \textbf{Fallback:} no certificate within budget \(\neq\) safety proof; candidate escalates to induction/CHC/PDR/concolic/triage.
  \item \textbf{Q\&A:} source mapping uses file/function/offset+symbol bindings; abstraction errors are caught by obligation counterexamples.
\end{itemize}
\normalsize
\end{frame}

\begin{frame}{7. Episode III: executable Python semantics}
\syncline{Slide 7}{~2 min}
\textbf{Script:}
\begin{itemize}
  \item Blog Episode III is the key turn: stop hand-waving and model CPython behavior explicitly.
  \item Include exception paths, call effects, and unknown behavior assumptions as first-class semantics.
  \item Representation split is intentional: concrete control, symbolic data/feasibility.
  \item Avoid overclaim: this is executable approximation with explicit assumptions, not full interpreter equivalence.
\end{itemize}
\end{frame}

\begin{frame}{8. Blog technical example: assertion escape}
\syncline{Slide 8}{~2 min}
\textbf{Script:}
\begin{itemize}
  \item Claim: ``failing assertion escapes uncaught'' is a reachability property, not a pattern match.
  \item Must account for assert mode, local handlers, caller handlers, and \texttt{finally} propagation.
  \item Naive rule-based detection over-reports; theorem-only abstractions under-specify behavior.
  \item The loop resolved this by combining executable semantics with conservative obligations.
  \item This is the blog's concrete argument for theory+code co-evolution.
\end{itemize}
\end{frame}

\begin{frame}{9. Transition system used by the analyzer}
\syncline{Slide 9}{~2 min}
\textbf{Script:}
\begin{itemize}
  \item State tuple: \(\sigma=(pc,\rho,h,\tau,\kappa,e,\Phi)\).
  \item \(pc\) and frame-stack control are concrete; branching uncertainty lives in \(\Phi\).
  \item \(\rho,h\) carry symbolic values and heap observers for solver-visible reasoning.
  \item Exception mode and handler context are explicit; this is critical for uncaught-escape correctness.
  \item Unsafe predicates are bug-family-indexed: \(U_t(\sigma)\).
\end{itemize}
\end{frame}

\begin{frame}{10. Real engine from the blog: theorize \(\rightarrow\) code \(\rightarrow\) test \(\rightarrow\) fix code \(\rightarrow\) fix theory}
\syncline{Slide 10}{~2 min}
\textbf{Script:}
\begin{itemize}
  \item This loop is not process fluff; it is the quality control mechanism.
  \item ``Fixing theory'' means tightening assumptions and obligations when tests expose semantic mismatch.
  \item Emphasize falsification pressure from synthetic suites and real-repo scans.
  \item Good soundbite: elegant ideas are cheap; executable invariants are expensive.
\end{itemize}
\end{frame}

\begin{frame}{11. What shipped: static-first, agentic-second}
\syncline{Slide 11}{~2 min}
\textbf{Script:}
\begin{itemize}
  \item Static stage handles the volume: discovery, feasibility, proofs, elimination, evidence packing.
  \item Agentic stage handles survivors only, with repository context exploration.
  \item Explicitly separate theorem-status artifacts from confidence/rationale artifacts.
  \item This asymmetry is both a trust design and a scaling design.
\end{itemize}
\end{frame}

\begin{frame}{12. Kitchensink portfolio: blog principle in formal-method terms}
\syncline{Slide 12}{~2 min}
\textbf{Script:}
\begin{itemize}
  \item Blog wording: ``steal the best ideas, orchestrate them, don't worship one paper.''
  \item Translation: route bug shapes through complementary proving families.
  \item Typical route: guard+SMT \(\rightarrow\) barrier attempts \(\rightarrow\) induction/CHC/PDR \(\rightarrow\) concolic confirmation.
  \item Keep disagreement visible; it is diagnostic signal for abstraction quality.
  \item Transition to next two slides: explain \textbf{why} this composition is coherent.
\end{itemize}
\end{frame}

\begin{frame}{12b. Barrier obligations as the unifying interface}
\syncline{Slide 12b}{~3 min}
\textbf{Script:}
\begin{itemize}
  \item The barrier certificate imposes three provable obligations: init, consecution, unsafe separation.
  \item These obligations are not just a proof artifact --- they are an \textbf{API}.
  \item Each paper from the 20-paper portfolio is a specialised subsolver for one or more of those obligations.
  \item The framework is not ``use everything and hope''; it is structured decomposition of a single verification problem.
  \item RISE point: the semantic contract (\(B\) with three conditions) is fixed; methods compete to discharge each condition.
  \item Method failure is informative: if SOS cannot certify separation, that tells you the unsafe region boundary is outside polynomial reach for the current variable basis.
  \item That failure is then a routing signal, not a dead end.
\end{itemize}
\end{frame}

\begin{frame}[shrink=62]{12c. Separation obligation: foundations (papers 1--5)}
\syncline{Slide 12c}{~4 min}
\textbf{Script:}
\begin{itemize}
  \item \textbf{BSP cluster role:} This cluster answers ``what does it mean for a barrier to exist and how do we compute one?'' Stengle guarantees existence; Parrilo makes it computable; Lasserre ensures degree-escalation; Prajna \& Jadbabaie wrote the three-obligation API we use verbatim.
  \item \textbf{Stengle 1974 (Positivstellensatz):} if \(\Phi(x)\wedge U(x)\) is infeasible over the reals, a polynomial identity certifies it. Purely theoretical existence guarantee under Archimedean conditions.\\
    \textbf{How we bent it:} the paper is pure real algebraic geometry with zero program semantics. We interpret program state variables as the polynomial ring and treat \texttt{assert} guards / error conditions as the semialgebraic sets \(S_0\) and \(U\).
  \item \textbf{Parrilo 2000 (SOS):} \(p\ge0\) certified by SOS decomposition \(p=\sum q_i^2\Leftrightarrow\) SDP over Gram matrix. Bridge from ``certificate exists'' to ``we can compute one.''\\
    \textbf{How we bent it:} Parrilo's domain is continuous nonlinear control (Lyapunov for ODEs). We substitute the ODE state space with RISE's symbolic machine state \(\sigma\), leaving the SDP machinery completely untouched.
  \item \textbf{Lasserre 2001 (moment/SOS hierarchy):} iterative degree-lifting converges to global optimum via moment sequences.\\
    \textbf{How we bent it:} framed as a polynomial \emph{optimisation} problem; we repurpose it as a \emph{feasibility-escalation} schedule where each hierarchy level is one tier of the portfolio, not an optimisation step.
  \item \textbf{Prajna \& Jadbabaie HSCC 2004:} first explicit three-obligation form. We use it verbatim over symbolic machine state.\\
    \textbf{How we bent it:} original uses Lie-derivative consecution (\(\dot B\ge0\) along ODE flow). We replace this with symbolic one-step \(B(\sigma')\ge0\) under the discrete transition relation \(T\) of the RISE IR.
  \item \textbf{Prajna, Jadbabaie \& Pappas TAC 2007:} consecution as expectation inequality; motivates stochastic risk path.\\
    \textbf{How we bent it:} replaces It\^{o} calculus over SDEs with Python probabilistic branching; expectation is approximated over path-condition probability mass drawn from the concrete execution distribution, not a stochastic differential equation.
  \item RISE point: the whole SOS chain (papers 1--5) establishes existence \(\rightarrow\) computability \(\rightarrow\) degree-escalation \(\rightarrow\) software interface.
\end{itemize}
\end{frame}

\begin{frame}[shrink=56]{12d. Separation tractability: relaxation ladder (papers 6--8)}
\syncline{Slide 12d}{~3 min}
\textbf{Script:}
\begin{itemize}
  \item \textbf{BSP cluster role:} Same fence, proved with cheaper paint. \(B\) still satisfies all three obligations; only the positivity proof is cheaper. DSOS/SDSOS are sound subsets of SOS: they never produce a false certificate.
  \item \textbf{Ahmadi \& Majumdar SIAM 2019 (DSOS/SDSOS):} DSOS replaces SDP with LP (diagonally-dominant SOS); SDSOS uses SOCP. Both are sound and 50--200\texttimes{} faster. First tier in the portfolio; escalate to full SDP only on failure.\\
    \textbf{How we bent it:} original targets robotics trajectory optimisation (LQR, control design). We discard the control objective entirely and use DSOS purely as a certificate-existence check.
  \item \textbf{Waki, Kim, Kojima \& Muramatsu 2006 (correlative sparsity):} decompose SDP along the variable correlation graph; reduces one large SDP to many small coupled SDPs.\\
    \textbf{How we bent it:} the original graph is built from the polynomial \emph{support} structure in the mathematical programme; we build it from the Python \emph{data-flow}: two variables co-appear iff they share a branch predicate or arithmetic expression in the RISE IR.
  \item \textbf{Ames et al.\ CBF survey 2019:} continuous consecution condition \(\dot B\ge -\alpha(B)\) (Lie derivative with class-\(\mathcal{K}\) bound).\\
    \textbf{How we bent it:} requires smooth continuous dynamics with computable Lie derivatives. We discretise the class-\(\mathcal{K}\) bound into a per-call contract \(B(\sigma')\ge B(\sigma)-\alpha(B(\sigma))\) and apply it to black-box Python callees, turning a continuous slope condition into a discrete overapproximation contract.
  \item RISE Q: ``Why not just SOS everywhere?'' DSOS is 50--200\texttimes{} faster; in CI that is 10 s vs.\ 30 min per scan.
\end{itemize}
\end{frame}

\begin{frame}[shrink=68]{12e. Consecution obligation: property-directed reachability (papers 9--12)}
\syncline{Slide 12e}{~4 min}
\textbf{Script:}
\begin{itemize}
  \item \textbf{BSP cluster role:} Handle consecution when no polynomial template is available. IC3/PDR build an implicit barrier as a clause conjunction; CHC encodes consecution as a horn-clause fixpoint. Same fence concept, represented symbolically.
  \item \textbf{Sheeran, Singh \& St\aa lm\aa rck FMCAD 2000 (k-induction):} splits consecution check into bounded base+step. After \(k\) unrolled steps without violation, promote candidate to bounded-safe.\\
    \textbf{How we bent it:} propositional SAT over hardware bit-vectors. We run SMT over integer/real program variables and upgrade the logic engine; base/step structure is adopted verbatim, then mapped onto our init- and consecution-obligation checks.
  \item \textbf{Bradley VMCAI 2011 (IC3):} frame predicates \(F_0\supseteq F_1\supseteq\ldots\) over-approximating reachable sets; blocking cubes until \(F_i=F_{i+1}\).\\
    \textbf{How we bent it:} Boolean SAT over hardware registers. We map ``cube'' to a conjunction of LIA predicates over RISE state, replace SAT-based cube generalisation with Z3 model-based quantifier instantiation, and accept IC3 as a heuristic (SAT completeness argument does not transfer to LIA).
  \item \textbf{Een, Mishchenko \& Brayton FMCAD 2011 (PDR):} industrial IC3 with ternary-simulation generalisation; gives UNSAT clause-sequence certificate.\\
    \textbf{How we bent it:} ternary simulation is hardware-specific. We replace it with SMT-driven predicate weakening. We \emph{do} retain PDR's proof-artifact extraction: the inductive clause sequence is serialised into the SARIF barrier-proof object.
  \item \textbf{Gurfinkel et al.\ CAV 2015 (SeaHorn/CHC):} software safety as CHC satisfiability over LLVM IR; head predicates = post-states, bodies = transitions.\\
    \textbf{How we bent it:} SeaHorn's CHC bodies are built from LLVM IR (C semantics with malloc/pointer). We re-encode for Python's RISE IR. Python dynamic dispatch requires per-callsite CHC heads rather than per-function, demanding a new CHC encoding layer that SeaHorn does not need.
  \item RISE point: IC3/PDR/CHC handle consecution when polynomial templates fail; they are the fallback tier.
\end{itemize}
\end{frame}

\begin{frame}[shrink=66]{12f. Template refinement: CEGAR and interpolants (papers 13--16)}
\syncline{Slide 12f}{~4 min}
\textbf{Script:}
\begin{itemize}
  \item \textbf{BSP cluster role:} Failed obligation counterexample = proof the current monomial basis cannot represent the fence; these papers extract the minimal new term needed to patch it.
  \item \textbf{Clarke, Grumberg, Jha, Lu \& Veith CAV 2000 (CEGAR):} abstract \(\rightarrow\) counterexample \(\rightarrow\) check concrete \(\rightarrow\) if spurious, refine. Classic termination argument via finite refinement steps.\\
    \textbf{How we bent it:} original ``abstraction'' is a boolean program; ours is the polynomial template degree and monomial basis. ``Refinement'' adds new monomials from the spurious trace path conditions instead of new boolean predicates.
  \item \textbf{Ball, Majumdar, Millstein \& Rajamani PLDI 2001 (SLAM):} CEGAR at Windows-driver scale; predicates from the counterexample trace seed the boolean abstraction.\\
    \textbf{How we bent it:} SLAM's predicate engine is tightly coupled to C boolean abstraction. We inherit only the \emph{architectural} lesson (automated refinement at scale) and substitute a polynomial monomial selector for SLAM's boolean predicate engine.
  \item \textbf{McMillan CAV 2003 (Craig interpolants):} UNSAT proof \(\Rightarrow\) interpolant formula implied by prefix, contradicting suffix. Strongest new predicate for refinement.\\
    \textbf{How we bent it:} original extracts propositional/linear interpolants for boolean abstraction. We apply SMT interpolation over non-linear arithmetic; the resulting formula is mapped to new monomial \emph{terms} in the polynomial template basis rather than boolean predicates.
  \item \textbf{Henzinger, Jhala, Majumdar \& Suresh POPL 2004 (lazy abstraction / BLAST):} refine only the failing subtree of the abstract reachability tree.\\
    \textbf{How we bent it:} ``subtree'' in BLAST is a CFG node set in boolean predicate abstraction. In our portfolio, it is the monomial support region of \(B\) around the failing obligation check; we localise the SDP re-solve to that variable sub-cluster, not the whole function.
  \item RISE point: these four papers together answer ``how do we improve the template without manual work?''
\end{itemize}
\end{frame}

\begin{frame}[shrink=73]{12g. Template synthesis and witness grounding (papers 17--20)}
\syncline{Slide 12g}{~4 min}
\textbf{Script:}
\begin{itemize}
  \item \textbf{BSP cluster role:} SOS says what \emph{shape} \(B\) must have; these papers find the concrete \(c_\alpha\) values and confirm the certificate faithfully reflects the real program. Closes the loop from Stengle (``a fence exists'') to ``we have the right fence for this specific bug site.''
  \item \textbf{Flanagan \& Leino FME 2001 (Houdini):} optimistically conjoin all candidate invariant clauses; drop those falsified by counterexample; iterate to stable fixpoint.\\
    \textbf{How we bent it:} Houdini proposes boolean invariant \emph{annotations} for ESC/Java method contracts. We generate polynomial \emph{coefficient} candidates instead; the pruning oracle is the three-obligation SMT check rather than an annotation type-checker.
  \item \textbf{Garg, L\"oding, Madhusudan \& Neider CAV 2014 (ICE):} learn invariants from positive / negative / implication examples; learner returns a classifier consistent with all three.\\
    \textbf{How we bent it:} ICE targets general first-order invariants with an unrestricted hypothesis class. We restrict to polynomials of bounded degree; the learner returns \(c_\alpha\) coefficient vectors rather than arbitrary formulas; concrete Python execution traces are the example source.
  \item \textbf{Alur, Bodik, Juniwal et al.\ FMCAD 2013 (SyGuS):} grammar-bounded synthesis with an SMT oracle; grammar constrains the candidate space, oracle checks correctness.\\
    \textbf{How we bent it:} SyGuS is entirely generic (bit manipulation, strings, etc.). We instantiate with a monomial grammar and supply all three BSP obligations as the multi-property SMT oracle, converting a general framework into a targeted certificate search.
  \item \textbf{Godefroid, Klarlund \& Sen PLDI 2005 (DART/concolic):} concrete+symbolic execution to maximise coverage and find assertion violations.\\
    \textbf{How we bent it:} DART's goal is \emph{falsification}. We \emph{invert} the objective: drive concolic runs toward certificate violations to \emph{confirm} \(B\ge0\) on all explored paths and \(B\le-\epsilon\) at real unsafe inputs, turning a bug-finder into a certificate validator.
  \item RISE summary: papers 17--20 close the loop. SOS = \emph{what shape}; Houdini/ICE/SyGuS = \emph{which coefficients}; concolic = \emph{is the certificate faithful?}
\end{itemize}
\end{frame}

\begin{frame}{13. RISE detail: solver outcomes and pruning policy}
\syncline{Slide 13}{~2 min}
\textbf{Script:}
\begin{itemize}
  \item \(\text{UNSAT}(\Phi\wedge U_t)\): eliminate with proof artifact.
  \item \(\text{SAT}(\Phi\wedge U_t)\): retain and possibly witness-seed.
  \item \(\text{UNKNOWN}/\text{timeout}\): retain conservatively; escalate stage.
  \item This policy is central to soundness posture under practical solver limits.
  \item If asked: yes, this can preserve extra noise; that is the deliberate safety tradeoff.
\end{itemize}
\end{frame}

\begin{frame}{14. Contracts for unknown/library calls}
\syncline{Slide 14}{~2 min}
\textbf{Script:}
\begin{itemize}
  \item Pure determinism for unknown calls is unsound; pure nondeterminism is too noisy.
  \item Middle ground: conservative contracts with widening and fallback assumptions.
  \item Contract precision is a first-order determinant of FP/TP balance.
  \item Concolic checks can confirm some SAT survivors when contracts are loose.
  \item State clearly: this is an explicit assumption interface, not hidden magic.
\end{itemize}
\end{frame}

\begin{frame}{15. CEGAR-style refinement}
\syncline{Slide 15}{~2 min}
\textbf{Script:}
\begin{itemize}
  \item Coarse abstraction first for tractable throughput.
  \item Spurious survivors trigger targeted refinement of predicates/contracts/context.
  \item Re-check obligations after refinement, not ad-hoc confidence tweaking.
  \item This is how precision improves without globally exploding compute.
\end{itemize}
\end{frame}

\begin{frame}{16. DSE/concolic as witness stage}
\syncline{Slide 16}{~2 min}
\textbf{Script:}
\begin{itemize}
  \item Concolic runs ground symbolic possibilities into executable witnesses when feasible.
  \item Witness presence increases actionability confidence for true positives.
  \item Witness absence is not a safety proof; bounded search limits apply.
  \item Keep hierarchy explicit: proof artifacts dominate confidence artifacts.
\end{itemize}
\end{frame}

\begin{frame}{17. Agentic-second triage}
\syncline{Slide 17}{~2 min}
\textbf{Script:}
\begin{itemize}
  \item Triage inspects source ranges, guards, callers, imports, tests, and nearby patterns.
  \item It classifies only survivors that static stage could not prove/disprove within budget.
  \item Emphasize governance: triage does not invalidate proven UNSAT eliminations.
  \item Better formal evidence upstream usually improves triage quality downstream.
\end{itemize}
\end{frame}

\begin{frame}{18. CI ratchet: not a firehose}
\syncline{Slide 18}{~2 min}
\textbf{Script:}
\begin{itemize}
  \item Baseline historical findings once; block new unaccepted risk in CI.
  \item Auto-prune resolved findings for monotonic quality movement.
  \item This process model is often the difference between adoption and abandonment.
  \item Phrase: ``ratchet turns infinite backlog into finite branch policy.''
\end{itemize}
\end{frame}

\begin{frame}{19. Soundness envelope and limits}
\syncline{Slide 19}{~2 min}
\textbf{Script:}
\begin{itemize}
  \item Dynamic Python features and reflective behavior bound static completeness.
  \item Unknown external behavior remains contract-limited.
  \item Resource budgets impose incompleteness on hardest nonlinear/interprocedural cases.
  \item Strong claim to make: assumptions are explicit and failure modes are observable.
\end{itemize}
\end{frame}

\begin{frame}{20. What RISE reviewers usually ask}
\syncline{Slide 20}{~2 min}
\textbf{Script:}
\begin{itemize}
  \item \textbf{State encoding:} what is concrete vs symbolic?
  \item \textbf{Solver policy:} how are UNKNOWN and timeout treated?
  \item \textbf{Contracts:} where do unknown-call assumptions live and how are they audited?
  \item \textbf{Refinement:} what triggers CEGAR and when do you stop?
  \item \textbf{Artifact chain:} can every elimination be explained post-hoc?
\end{itemize}
\end{frame}

\begin{frame}[fragile]{21. Operational run modes}
\syncline{Slide 21}{~2 min}
\textbf{Script:}
\begin{itemize}
  \item Quick demo mode: \texttt{a3 scan . --triage}
  \item Explicit CI mode: \texttt{a3 scan . --output-sarif results.sarif}, then \texttt{a3 triage --sarif results.sarif --agentic}
  \item Keep stage separation for debugging, traceability, and policy governance.
  \item For reproducibility in CI, pin model/provider and archive SARIF artifacts.
\end{itemize}
\end{frame}

\begin{frame}{22. Close + Q\&A}
\syncline{Slide 22}{~2 min}
\textbf{Script:}
\begin{itemize}
  \item Re-state three claims: static-first trust, conservative escalation, ratchet adoption.
  \item Invite RISE discussion on soundness envelope, obligation encoding, and portfolio scheduling.
  \item Suggested starting point in practice: run static stage with SARIF evidence and baseline ratchet before enabling triage.
  \item Future work line: stronger contracts and richer nonlinear invariant synthesis under strict budgets.
\end{itemize}
\end{frame}

\end{document}
